\chapter{Frontiers}\label{ch:frontiers}

Open problems across all the parts, and a guide to the recent literature.

\section{Open problems across all parts}\label{ch:open_all}

\begin{physmotiv}
A map of what is not known, organized by the Parts of this book.
\end{physmotiv}

\begin{description}[leftmargin=1em]
\item[Parts II--III (algebraic foundations).] Classification of factors beyond the hyperfinite case is wide open; which type III$_1$ factors arise as local algebras of interacting QFTs in $d=4$? Does the hyperfinite property (injectivity) of local algebras hold in general, or only under nuclearity assumptions?
\item[Part IV (modular theory).] Explicit modular Hamiltonians for generic regions and states in interacting QFT are known only in special cases (wedges, balls in CFT, a few free-field examples). Is there a geometric characterization of when the modular flow is local?
\item[Part VI (AQFT).] Existence of interacting nets in $3+1$ dimensions; the relation between algebraic and constructive QFT.
\item[Part VII (de Sitter and observers).] See \cref{ch:open_ds}: non-perturbative definition, observer dependence, the additive constant, multi-observer algebras, dynamics.
\item[Part VIII (black holes).] The finite-$N$ origin of the large-$N$ type III$_1$ limit; the algebraic description of black hole evaporation including the island formula; the type of the algebra of the exterior of an evaporating black hole.
\item[Part IX (noncommutative geometry).] A Lorentzian spectral triple formalism; the status of the spectral Standard Model; the relationship, if any, to the gravitational algebras of Part VII.
\end{description}

\begin{keyidea}
The unifying open question: \emph{what is the algebra of observables of quantum gravity, and which of its properties (type, trace, entropy) are robust beyond the semiclassical limit?}
\end{keyidea}

\section{Guide to the recent literature}\label{ch:literature}

\begin{physmotiv}
A reading guide, grouped by topic. Entries marked as verified were checked against their arXiv/journal records while writing; others are recalled and need verification before reliance (see \texttt{refs.bib}).
\end{physmotiv}

\subsection{Textbooks}
Bratteli--Robinson \cite{BratteliRobinson1,BratteliRobinson2} (operator algebras and quantum statistical mechanics; KMS, quasi-local algebras); Haag \cite{Haag} (local quantum physics); Takesaki \cite{Takesaki1,Takesaki2} (the standard mathematical reference for modular theory and crossed products); Kadison--Ringrose \cite{KadisonRingrose1}; Reed--Simon \cite{ReedSimon1} for functional analysis; Connes \cite{Connes1994} for noncommutative geometry.

\subsection{Entanglement in QFT and gravity: algebraic perspective}
Witten's review of entanglement properties of QFT \cite{Witten2018}; Witten, ``Gravity and the crossed product'' \cite{Witten2021}; Chandrasekaran--Penington--Witten \cite{CPW2022}; Chandrasekaran--Longo--Penington--Witten (de Sitter) \cite{CLPW2023} (verified).

\subsection{Generalized entropy for subregions and observers}
Jensen--Sorce--Speranza \cite{JSS2023} (verified); Faulkner--Speranza \cite{FaulknerSperanza2024} (verified); Chen--Penington \cite{ChenPenington2024} (verified); De Vuyst--Eccles--H\"ohn--Kirklin \cite{DEHK2024a,DEHK2024b} (verified); Kirklin \cite{Kirklin2024GSL} (title verified; details to check); Witten's background-independent algebra \cite{Witten2023bg} (not verified).

\subsection{Foundations of de Sitter thermodynamics}
Gibbons--Hawking \cite{GibbonsHawking1977}. Connes--Rovelli on thermal time \cite{ConnesRovelli1994}.

\subsection{Spectral geometry}
Chamseddine--Connes \cite{CCM1997}; Connes' reconstruction theorem (2008, 2013); the Chamseddine--Connes--Marcolli paper on gauge theory and the Standard Model; the recent literature on the spectral Standard Model (to be checked).

\subsection{How to keep up}
The papers in this subject appear mostly in hep-th and math.OA; search by author and by the keywords ``crossed product'', ``type II$_1$'', ``generalized entropy'', ``quantum reference frame''. Since this literature is moving, treat the summaries in \cref{ch:clocks_qrf,ch:gen_entropy} as a snapshot as of late 2026 and verify before citing.

